\def\probtitle{개성있는 구간}
\def\probno{PB} % 문제 번호

\begin{problem}{\probtitle}{표준 입력(stdin)}{표준 출력(stdout)}{1 초}{1024 megabytes}

다음 조건을 만족하는 길이가 $N$인 수열 $A = [A_1, A_2, \cdots, A_N]$를 아무거나 하나 찾아라.

\begin{itemize}[topsep=0pt,noitemsep]
    \item 수열 $A$의 모든 원소는 $1$ 이상 $6N$ 이하이다.
    \item $1 \le L \le R \le N$을 만족하는 모든 정수 쌍 $(L, R)$에 대하여, $\sum\limits_{i=L}^R A_i$가 가질 수 있는 서로 다른 값은 모두 다르다.
\end{itemize}

주어진 범위 내에서 조건을 만족하는 수열은 항상 존재한다.

\InputFile

첫째 줄에 수열 $A$의 길이를 나타내는 정수 $N$이 주어진다.

\OutputFile

첫째 줄에 조건을 만족하는 수열 $A$의 원소들을 공백으로 구분하여 출력한다.

위 조건을 만족하는 수열이 여러 개라면 그중 아무거나 출력한다.

\InputFile

첫째 줄에 정수 $N$ 이 주어진다.

\OutputFile

첫째 줄에 정수 $N$개 $A_1, A_2, \dots, A_N$을 공백으로 구분하여 출력하라.

\Constraints

\begin{itemize}[topsep=0pt,noitemsep]
    \item $1 \le N \le 8192$
    \item $1 \le i \le N$인 모든 $i$ 에 대해 $1 \le A_i \le 6N$
\end{itemize}

\Example

\begin{example}
    \exmpfile{./example/01.in.txt}{./example/01.out.txt}%
\end{example}

\Notes

예제 출력의 수열은 $A = [2, 4, 8, 16]$이다.

이때 모든 연속 부분수열의 합은 다음과 같다.

\begin{itemize}[topsep=0pt,noitemsep]
    \item $[1, 1]$: $2$
    \item $[2, 2]$: $4$
    \item $[3, 3]$: $8$
    \item $[4, 4]$: $16$
    \item $[1, 2]$: $2 + 4 = 6$
    \item $[2, 3]$: $4 + 8 = 12$
    \item $[3, 4]$: $8 + 16 = 24$
    \item $[1, 3]$: $2 + 4 + 8 = 14$
    \item $[2, 4]$: $4 + 8 + 16 = 28$
    \item $[1, 4]$: $2 + 4 + 8 + 16 = 30$
\end{itemize}

따라서 모든 연속 부분수열의 합은
$2, 4, 8, 16, 6, 12, 24, 14, 28, 30$으로 서로 다르다.

\end{problem}